% !TeX root = ../../../main.tex
\section{The common group-scheme construction}
\label{appendix:common_group_scheme_construction}

The three cyclic theories used in Chapter~2 are instances of the following
construction.

Let $H$ be a commutative group scheme over a scheme $S$, written additively,
and let $\varphi:H\to H$ be a homomorphism. Its scheme-theoretic kernel
$K:=\ker(\varphi)$ is defined by the Cartesian square
\begin{equation}\label{eq:kernel_square}
\begin{tikzcd}
K \arrow[r] \arrow[d] & H \arrow[d,"\varphi"] \\
S \arrow[r,"0_H"'] & H.
\end{tikzcd}
\end{equation}
Thus, for every $S$-scheme $U$,
\begin{equation}\label{eq:kernel_on_U_valued_points}
K(U)=\ker(\varphi_U)
    =\{k\in H(U):\varphi_U(k)=0_{H(U)}\},
\end{equation}
where $H(U)=\Hom_S(U,H)$ and $\varphi_U$ is induced by $\varphi$.

For an $S$-scheme $X$ and a section $a\in H(X)$, define the fiber
\begin{equation}\label{eq:fiber_torsor_definition}
T_a:=X\times_{a,H,\varphi}H,
\end{equation}
equivalently by the Cartesian square
\[
\begin{tikzcd}
T_a \arrow[r] \arrow[d] & H \arrow[d,"\varphi"] \\
X \arrow[r,"a"'] & H.
\end{tikzcd}
\]
For every $X$-scheme $U\to X$, with $a_U$ the pullback of $a$,
\[
T_a(U)=\{y\in H(U):\varphi_U(y)=a_U\}.
\]
Thus ``the torsor $\varphi(Y)=a$'' means the scheme of solutions after every
base change to $X$; it has no preferred identity element in general.

Write $H_X:=H\times_S X$, $K_X:=K\times_S X$, let
$\varphi_X:H_X\to H_X$ be the base change of $\varphi$, and let
$\iota_K:K_X\to H_X$ and $\iota_a:T_a\to H_X$ be the canonical maps.
Addition on $H_X$ restricts to the action
\begin{equation}\label{eq:kernel_action_on_fiber}
K_X\times_X T_a\longrightarrow T_a,
\qquad (k,y)\longmapsto y+k.
\end{equation}
It is the top arrow in the commutative diagram
\[
\begin{tikzcd}[column sep=large, row sep=large]
K_X\times_X T_a
    \arrow[r,"{(k,y)\mapsto y+k}"]
    \arrow[d,"{\iota_K\times_X\iota_a}"']
  & T_a \arrow[d,"\iota_a"] \\
H_X\times_X H_X
    \arrow[r,"+_X"']
  & H_X.
\end{tikzcd}
\]
Indeed,
\[
\varphi_U(y+k)
=\varphi_U(y)+\varphi_U(k)
=a_U+0_{H(U)}
=a_U,
\]
so translation by $K$ preserves the fiber $T_a$. The action is simply
transitive: the morphism
\begin{equation}\label{eq:torsor_difference_isomorphism}
K_X\times_X T_a\longrightarrow T_a\times_X T_a,
\qquad (k,y)\longmapsto (y,y+k),
\end{equation}
is an isomorphism with inverse
\[
(y_1,y_2)\longmapsto (y_2-y_1,y_1).
\]
This inverse is well-defined because
\[
\varphi_U(y_2-y_1)
=\varphi_U(y_2)-\varphi_U(y_1)
=a_U-a_U
=0_{H(U)}.
\]
If $\varphi$ is finite \'etale and surjective, then its base change
$T_a\to X$ is finite \'etale and surjective. Hence $T_a$ is a $K_X$-torsor,
namely
\[
T_a=a^*(H\xrightarrow{\varphi}H)
\]
as a pullback of the $K$-torsor $\varphi:H\to H$.

Suppose that $a'=a+\varphi_X(b)$ for some $b\in H(X)$. Translation by $b$
defines
\[
\tau_b:T_a\longrightarrow T_{a'},
\qquad y\longmapsto y+b,
\]
because
\[
\varphi_U(y+b_U)
=a_U+\varphi_U(b_U)
=a'_U,
\]
and it is $K_X$-equivariant since
$\tau_b(y+k)=y+k+b=\tau_b(y)+k$. Its inverse is $\tau_{-b}$, so $T_a$ depends
only on the class of $a$ modulo $\varphi_X(H(X))$.

When $\varphi$ is finite \'etale and surjective, regard the group schemes
$K_X$ and $H_X$ as the sheaves that they represent on the \'etale site of
$X$. Then
\[
0\longrightarrow K_X\longrightarrow H_X
\xrightarrow{\varphi_X}H_X\longrightarrow0
\]
is a short exact sequence of sheaves. Applying \'etale cohomology gives the
beginning of the associated long exact sequence
\[
0\longrightarrow H^0_{\mathrm{\acute et}}(X,K_X)
\longrightarrow H^0_{\mathrm{\acute et}}(X,H_X)
\xrightarrow{\varphi_X}H^0_{\mathrm{\acute et}}(X,H_X)
\xrightarrow{\delta}H^1_{\mathrm{\acute et}}(X,K_X).
\]
Because $H_X$ is representable, its global sections are
\[
H^0_{\mathrm{\acute et}}(X,H_X)
=\Hom_X(X,H_X)
=\Hom_S(X,H)
=H(X).
\]
For $a\in H(X)$, the connecting class $\delta(a)$ is represented by the
sheaf of local lifts of $a$ through $\varphi_X$. This sheaf is precisely the
pullback $T_a=X\times_{a,H,\varphi}H$, with its $K_X$-action, and hence
$\delta(a)=[T_a]$. Exactness gives
$\ker(\delta)=\varphi_X(H(X))$, so $\delta$ descends to the injection
\begin{equation}\label{eq:common_torsor_connecting_map}
H(X)/\varphi_X(H(X))\lhook\joinrel\longrightarrow
H^1_{\mathrm{\acute et}}(X,K_X),
\qquad [a]\longmapsto[T_a].
\end{equation}
The class $[T_a]$ is trivial exactly when $a\in\varphi_X(H(X))$, equivalently
when $\varphi(Y)=a$ has a global solution.

For a multiplicative group scheme the same construction is written
multiplicatively: $K(U)=\{k\in H(U):\varphi_U(k)=1_{H(U)}\}$, the action is
$(k,y)\mapsto ky$, and replacing $a$ by $a\varphi_X(b)$ gives the isomorphic
torsor $y\mapsto by$.

\subsection*{The three cyclic theories}

For an arbitrary $S$-scheme $X$, the fiber over a section $a\in H(X)$ is
\[
T_a=\varphi_X^{-1}(a)
   =X\times_{a,H,\varphi}H.
\]
For nonaffine $X$, explicit equations are interpreted using the structure
sheaf $\Oo_X$. In all three applications in Chapter~2, however, the base is
affine: first $\Spec k((x))$ and then
$\bG_m=\Spec k[x,x^{-1}]$. We therefore work affinely from now on and write
\[
X=\Spec A,
\]
where the capital letter $A$ is the coordinate ring of $X$, while the
lowercase letter $a$ is the section specifying the fiber. More generally, if
$S=\Spec R$ and $H=\Spec C$, then
\[
T_a=\Spec\bigl(A\otimes_C C\bigr),
\]
where $C$ acts on $A$ through $a^\#:C\to A$ and on the second copy of $C$
through $\varphi^\#:C\to C$.

\paragraph{Kummer.}
Take $H=\bG_m$ and $\varphi(z)=z^e$. A section is an element
$a\in\bG_m(X)=A^\times$, and its fiber is
\[
T_a=\Spec\bigl(A[Y]/(Y^e-a)\bigr),
\]
where $Y$ is automatically a unit because $a$ is a unit. The kernel is
\[
\mu_{e,X}=\Spec\bigl(A[t,t^{-1}]/(t^e-1)\bigr),
\]
and it acts on $T_a$ through a morphism
\[
\alpha:\mu_{e,X}\times_X T_a\longrightarrow T_a.
\]
Writing
\[
B=A[Y]/(Y^e-a)
\quad\text{and}\quad
D=A[t,t^{-1}]/(t^e-1),
\]
the corresponding homomorphism of coordinate rings is
\[
\alpha^\#:B\longrightarrow D\otimes_A B,
\qquad
\overline Y\longmapsto t\otimes\overline Y,
\]
where $\overline Y$ denotes the class of $Y$ in $B$. This is well-defined
because
\[
(t\otimes\overline Y)^e=t^e\otimes\overline Y^e=1\otimes a.
\]
Equivalently, for every $X$-scheme $U$,
\[
T_a(U)=\{y\in\Gamma(U,\Oo_U)^\times:y^e=a_U\},
\qquad
\mu_e(U)=\{\xi\in\Gamma(U,\Oo_U)^\times:\xi^e=1\},
\]
and the action sends $(\xi,y)$ to $\xi y$; indeed,
$(\xi y)^e=\xi^e y^e=a_U$. If $e$ is invertible on $S$ and a primitive
root $\zeta_e\in\mu_e(S)$ is chosen, then
$\overline j\mapsto\zeta_e^j$ identifies $(\Z/e\Z)_S$ with $\mu_e$.
Only this identification, not the $\mu_e$-torsor or its action, depends on
the choice of $\zeta_e$.

\paragraph{Artin--Schreier.}
Let $A$ be an $\F_p$-algebra, take $H=\bG_a$, and put
$\varphi(Y)=Y^p-Y$. A section is an element $a\in\bG_a(X)=A$, and its
fiber is
\[
T_a=\Spec\bigl(A[Y]/(Y^p-Y-a)\bigr).
\]
The kernel is the constant group scheme $(\F_p)_X$, acting by
$(c,Y)\mapsto Y+c$.

\paragraph{Artin--Schreier--Witt.}
Let $A$ be an $\F_p$-algebra. The third theory takes $H=\bW_r$ and
\[
\varphi=\wp:=F-\mathrm{id},
\qquad
\wp\vec Y=F\vec Y\ominus\vec Y.
\]
For a section $\vec a\in W_r(A)$, the fiber is explicitly
\[
T_{\vec a}
=\Spec\!\left(
  A[Y_0,\dots,Y_{r-1}]
  \big/\bigl(\wp\vec Y\ominus\vec a\bigr)
\right),
\qquad \wp\vec Y=\vec a.
\]
Here the parenthesized Witt vector denotes the ideal generated by its $r$
coordinate polynomials. Its kernel is
$W_r(\F_p)_X\cong(\Z/p^r\Z)_X$, acting by Witt translation
$(\vec c,\vec Y)\mapsto\vec Y\oplus\vec c$.

This paragraph states the affine model in advance for comparison with the
two preceding theories. The operations $F$, $\oplus$, and $\ominus$, and the
coordinate polynomials occurring in the displayed quotient, are defined in
\Cref{appendix:witt_vector_algebra_basics}. The fact that this fiber is
finite \'etale and that the displayed kernel action makes it a torsor is
proved in \Cref{lemma:asw_etale_torsor}.

